\chapter{Evaluación}
\label{cap:evaluacion}

Este capítulo presenta los resultados empíricos del sistema SGROAS organizados por
bloque de evaluación (Bloque C de la guía \cite{uteq2026guia}). Cada bloque declara:
la pregunta de investigación que responde, la métrica empleada, los datos crudos que
la soportan, la estadística descriptiva con intervalo de confianza al 95\,\% y la
interpretación. Todos los números son generados con scripts versionados y pueden
reproducirse siguiendo \texttt{docs/mediciones/DATA-PROVENANCE.md}. El diseño
(GQM, RQs y umbrales declarados \emph{a priori}) está en
\texttt{docs/estudio/DISENO-ESTUDIO.md}.

\section{Rendimiento (Bloque C.1) — RQ1}
\label{sec:perf}

\textbf{Pregunta:} \emph{¿El acceso híbrido a datos CRUD/ORM + stored procedures
mantiene un tiempo de respuesta aceptable bajo carga?}

\textbf{Métricas:} latencia de \texttt{http\_req\_duration} (media, mediana,
percentil 95) y tasa de error \texttt{http\_req\_failed} sobre
\texttt{GET /api/conductores}, con 50 VUs durante 30\,s (\texttt{k6/opts.js}).
El umbral declarado es p95~<~200\,ms y error~<~1\,\%.

\textbf{Datos crudos:} tres corridas independientes en
\texttt{docs/mediciones/perf/k01-run1.json}, \texttt{k02-run2.json} y
\texttt{k03-run3.json} (commit \texttt{62bf8fa}), más cinco corridas contra la URL
pública en \texttt{k04-run1.json}--\texttt{k08-run1.json} y cinco muestras frías en
\texttt{k04-cold.json}--\texttt{k08-cold.json} (commit \texttt{9ded2a69}).

\textbf{Descriptiva e IC 95\,\%:} la media de las medias por corrida es
\textbf{23.01\,ms} (DT 12.44\,ms, EE 7.18\,ms), con IC\,95\,\%
\textbf{[$-7.88$; $53.91$]\,ms} (t de Student, gl = 2, n = 3 corridas). El p95
máximo es 173.13\,ms (corrida k01), por debajo del umbral de 200\,ms, y la tasa de
error es 0\,\% con 8930 verificaciones correctas. La primera corrida (37.3\,ms de
media; p95 173.13\,ms) refleja el arranque en frío de conexiones (pool JDBC,
JIT); las corridas k02 y k03 estabilizan la media por debajo de 17\,ms
(\texttt{docs/mediciones/perf/ANALISIS-k6.md}). Esta serie (K1--K3, commit
\texttt{62bf8fa}, 2026-07-29) es anterior a la implementación real de
\texttt{@Cacheable} en este endpoint (commit \texttt{3f78036}, 2026-09-17):
\texttt{GET /api/conductores} no tenía caché de aplicación activa en ese
momento, así que la caída de latencia entre corridas no puede atribuirse a
Redis, solo al arranque en frío del propio proceso.

\textbf{Interpretación:} el sistema cumple el umbral de p95~<~200\,ms en el
entorno local sin throttling (serie K1--K3), aunque esta serie no mide un
efecto de caché real -- ver el párrafo anterior. El contraste formal de
caché caliente/fría, con metodología válida, se ejecutó después contra la
URL pública una vez implementado \texttt{@Cacheable} (K9 y K10--K14, más
abajo).

\textbf{Serie Render (K4--K8, URL pública):} con la URL pública activa se ejecutaron
cinco corridas adicionales contra \texttt{https://sgroas-backend.onrender.com}
(\texttt{docs/mediciones/perf/k04-run1.json}--\texttt{k08-run1.json}), cinco muestras
frías (\texttt{k04-cold.json}--\texttt{k08-cold.json}) y el contraste no paramétrico
(\texttt{RENDER-REPORT.md}). Estas cinco corridas (50 VUs) se ejecutaron el
6 de septiembre, antes de corregirse un bug de auto-invocación de Spring AOP
que impedía que \texttt{@Cacheable} se activara en ningún servicio (ver
capítulo de implementación): \emph{ninguna} de las dos condiciones mide
caché real, así que aquí se llaman "con Redis reiniciado" y "con Redis
activo" -- no "fría" y "caliente" -- para no sugerir un efecto que esta
serie no puede medir. La media de avg con Redis activo es
\textbf{3967.32\,ms} (DT 1539.20, IC\,95\,\% [2056.46; 5878.19]) y con
Redis reiniciado, \textbf{257.50\,ms}. El p95 con Redis activo va de
4.3\,s a 11.4\,s; el plan gratuito de Render (0.1\,vCPU) no cumple el
umbral $p95 < 200$\,ms. La diferencia de orden de magnitud corresponde al
throttling de CPU del plan gratuito y al almacenamiento local de Render
\cite{jiang2015}. El contraste Mann-Whitney entre ambas condiciones arroja
$U = 0.0$, $p = 0.009$ y $d$ de Cliff $= -1.00$ (efecto grande), pero mide
sobre todo el arranque en frío de la instancia gratuita de Render, no el
efecto de la caché de aplicación -- que en este momento del proyecto ni
siquiera estaba activa. La serie local K1 (p95 medio 81\,ms) sigue siendo
la referencia de rendimiento sin throttling.

\textbf{Contraste frío/caliente de caché, con metodología corregida
(K9):} el contraste anterior comparaba perfiles de carga distintos (1
usuario/1 iteración frío vs.\ 50 usuarios/30\,s caliente), una comparación
no válida señalada en una auditoría externa. Corregido con
\texttt{k6/cache-contrast.js}: 1 solo usuario, misma corrida secuencial,
contra la URL pública ya con el fix de caché aplicado. La primera petición
(clave de caché nunca solicitada antes) es el \emph{miss} real; las diez
siguientes, contra la misma clave, son \emph{hits} reales
(\texttt{docs/mediciones/perf/k09-cache-contrast.json}, 2026-09-18):

\begin{center}
\begin{tabular}{@{}lccc@{}}
\toprule
Condición & Mediana & p95 & Máx \\
\midrule
Frío (miss, n=1) & 514.75\,ms & 514.75\,ms & 514.75\,ms \\
Caliente (hit, n=10) & 191.59\,ms & 504.77\,ms & 721.00\,ms \\
\bottomrule
\end{tabular}
\end{center}

La caché reduce la latencia típica (mediana) en un 63\,\% (514.75\,ms
$\rightarrow$ 191.59\,ms), un efecto real y medido con el mismo usuario y el
mismo perfil de carga en ambas condiciones. Sin embargo, el p95 en caliente
(504.77\,ms) sigue sin cumplir el umbral de 200\,ms: en el plan gratuito de
Render la variabilidad de red/CPU domina sobre el ahorro de consultar Redis
en vez de PostgreSQL, así que la caché no basta por sí sola para cumplir el
umbral de rendimiento en ese entorno. El umbral sí se cumple en el entorno
local sin throttling (serie K1).

\textbf{Contraste con 5 corridas (K10--K14), análisis no paramétrico:}
K9 fue una sola corrida piloto ($n=1$ en frío); la guía exige cinco
corridas por escenario analizadas con métodos no paramétricos. Se
repitió la misma metodología (mismo usuario, misma corrida secuencial,
\texttt{k6/cache-contrast.js}) cinco veces contra la URL pública, cada
una con una clave de caché distinta (\texttt{page size} 21 a 34) para
garantizar un \emph{miss} real e independiente en cada corrida. Estas
cinco son exportaciones \texttt{--summary-export} reales de k6, sin
editar (\texttt{docs/mediciones/perf/k10}--\texttt{k14-cache-contrast.json},
solo con el token de sesión de \texttt{setup\_data} redactado antes de
versionar). Recalculado desde los JSON crudos con
\texttt{scripts/perf/recalcular-contraste-cache.py}
(Mann-Whitney + $d$ de Cliff, \texttt{scripts/perf/nonparametric.py}):

\begin{center}
\begin{tabular}{@{}lccccc@{}}
\toprule
Corrida & K10 & K11 & K12 & K13 & K14 \\
\midrule
Frío (avg, ms) & 220.58 & 203.39 & 205.00 & 319.50 & 225.76 \\
Caliente (avg, ms) & 206.19 & 239.71 & 211.21 & 223.66 & 258.02 \\
\bottomrule
\end{tabular}
\end{center}

$U = 10{,}0$, $z = -0{,}52$, $p = 0{,}6015$ (bilateral), $d$ de Cliff $=
-0{,}20$ (efecto pequeño). Con datos ya genuinamente crudos (no la
reconstrucción de una versión anterior de este documento, corregida
tras una re-evaluación externa que señaló, con razón, que esos JSON no
tenían la estructura de una exportación real de k6) el resultado es
honesto y menos favorable de lo que parecía: no hay diferencia
significativa entre frío y caliente. De las cinco corridas, la
condición caliente resultó igual o más lenta que la fría en
\textbf{tres de cinco} (K11: 239{,}71 vs.\ 203{,}39; K12: 211{,}21 vs.\ 205{,}00; K14:
258{,}02 vs.\ 225{,}76) y más rápida solo en dos (K10, K13). Si la caché
redujera la latencia de forma consistente, la condición caliente
debería ganar en casi todas las corridas, no en 2 de 5.

La razón más plausible no es que la caché esté rota, sino que su efecto
real es demasiado pequeño para distinguirlo del ruido de la propia
medición: todas las diferencias observadas son de apenas 15--35\,ms,
mientras que la variabilidad de red/CPU del plan gratuito de Render
(ya documentada más arriba, con picos de segundos completos bajo carga)
es un orden de magnitud mayor. Es análogo a intentar medir si alguien
camina más rápido con zapatos nuevos usando una cinta métrica que tiembla
$\pm 2$ metros: si el efecto real es de unos centímetros, el temblor del
instrumento lo esconde por completo, y las lecturas terminan pareciendo
aleatorias. Con $n=5$ por condición no se puede distinguir el efecto de
caché del ruido de red. La única corrida que sí mostró una diferencia
grande (K9, frío 514{,}75\,ms vs.\ caliente mediana 191{,}59\,ms) sigue
siendo real, pero no representativa de las cinco corridas siguientes:
la caché no está rota (los tiempos calientes rara vez superan los
300\,ms, salvo el ruido de red ya señalado), pero esta medición no
alcanza a demostrar estadísticamente que la reduzca de forma
consistente en este entorno.

\section{Seguridad (Bloque C.2) — RQ2}
\label{sec:sec}

\textbf{Pregunta:} \emph{¿El sistema mitiga los riesgos más críticos del OWASP Top
10 \cite{owasp2021} con controles verificables?}

\textbf{Evidencia (datos crudos):} se documentan seis controles en
\texttt{docs/mediciones/sec/}: control de acceso (A01), criptografía de contraseñas
y sesión (A02), prevención de inyección con repositorios Spring Data y
\texttt{@Procedure} sin SQL concatenado (A03), cabeceras de seguridad y error
\emph{Problem Details} RFC\,7807 \cite{rfc7807} (A05), limitación de tasa de login
(A07) y registros de auditoría (A09).

\textbf{Medidas mecánicas:} la regla de oro del proyecto prohíbe SQL dinámico
concatenado; se audita con \texttt{scripts/audit-sql-dynamic.sh} (commit
\texttt{1a07dc7}) y con SpotBugs/Find-Sec-Bugs en CI. El esquema de autenticación
usa JWT según RFC\,7519 \cite{rfc7519} sobre OAuth\,2.0 \cite{rfc6749}, con cookie
HttpOnly y BCrypt (factor 10+).

\textbf{ZAP baseline (K3):} el escaneo baseline de OWASP ZAP
(\texttt{scripts/zap/run-zap.sh}) se ejecutó contra la URL pública
\texttt{https://sgroas-backend.onrender.com} (commit \texttt{ca7f700}). Resultado: 0
High, 2 Medium (configuración CSP, no vulnerabilidad), 4 Low (cabeceras CORS del
reverse proxy) e 2 Informational; 63 PASS sobre 8 URLs
(\texttt{docs/mediciones/sec/zap/RESUMEN.md}).

\textbf{Interpretación:} los controles estáticos y dinámicos ejecutados no detectan
SQL dinámico ni credenciales en texto plano; la evidencia A01 (accessible via
token) confirma el cierre del acceso no autenticado.

\textbf{Verificación en vivo (sesión y cookie segura):} contra el deploy público
(\texttt{https://sgroas-backend.onrender.com}) se verifica el flujo completo de
sesión: el login responde \texttt{200} con dos cookies
\texttt{Secure; HttpOnly; SameSite=Strict} (access 1\,h, refresh 7\,días), la petición
autenticada \texttt{GET /api/auth/me} responde \texttt{200} (ROLE\_ADMIN) y el
endpoint \texttt{GET /api/asignaciones?page=0\&size=10} con la cookie responde
\texttt{200} (8 asignaciones); sin cookie el acceso es rechazado (\texttt{403} en el
proxy). Evidencia cruda y comandos reproducibles en
\texttt{docs/mediciones/sec/live-session/REPORT.md}.

\section{Usabilidad (Bloque C.3) — RQ3}
\label{sec:sus}

\textbf{Pregunta:} \emph{¿Qué nivel de usabilidad perciben los usuarios del
sistema?}

\textbf{Instrumento:} System Usability Scale de Brooke \cite{brooke1996}, 10 ítems
Likert 1--5. Puntuación = (suma de contribuciones) $\times$ 2.5.

\textbf{Datos crudos:} 10 participantes (P01--P10) anonimizados en
\texttt{dataset/sus/sus-raw.csv} y por participante en
\texttt{P01.json..P10.json}; consentimientos custodiados fuera del repositorio.
Se recolectaron respuestas adicionales de 5 personas más (P11--P15), con
consentimiento real y verificado, pero sin respaldo documental verificable de
la fecha de recolección de sus respuestas al cuestionario; se excluyen del
análisis oficial (ver \texttt{dataset/sus/PARTICIPANTES-NO-INCLUIDOS.md}).

\textbf{Descriptiva e IC 95\,\%:} media \textbf{63,0/100}, DT 13,88, EE 4,39, IC\,95\,\%
\textbf{[53,07; 72,93]} (t = 2,262, gl = 9), mínimo 47,5 y máximo 90,0
(\texttt{dataset/sus/REPORT.md}). En la escala adjetiva de Bangor et
al. \cite{bangor2009,bangor2008} la media corresponde a \emph{OK}, dentro de la
zona de aceptabilidad \emph{marginal} (50--70) \cite{lewis2009}. La
calificación \emph{Bueno} de Bangor empieza alrededor de 70--72, por
encima del valor observado.

\textbf{Interpretación:} la media se ubica por debajo del umbral de 70 que Bangor
asocia a sistemas aceptables, por lo que el resultado es \emph{marginal} y se
documenta como área de mejora. El tamaño muestral (n = 10) se ubica en el rango
recomendado para SUS \cite{kortum2013}.

\section{Cobertura de pruebas (Bloque C.4)}
\label{sec:cobertura}

\textbf{Métrica:} contadores de JaCoCo (instrucciones, ramas y líneas) sobre la
ejecución de \texttt{./mvnw verify} con 299 pruebas JUnit 5.

\textbf{Datos crudos:} \texttt{docs/mediciones/jacoco/jacoco.csv} y reporte HTML
(versionado en \texttt{dataset/jacoco}).

\textbf{Resultado (reporte completo):} instrucciones 95.67\,\%,
ramas 88.38\,\% y líneas 96.19\,\%, con 0 fallos. Se supera el umbral del proyecto
(líneas $\ge$ 0.70 y ramas $\ge$ 0.50 para la entrega final). La discusión sobre la relación entre cobertura y detección de fallos se
fundamenta en \cite{inozemtseva2014,gopinath2014,gligoric2013,ivankovic2019}.

\section{Calidad web (Bloque C.5) — RQ4}
\label{sec:lighthouse}

\textbf{Pregunta:} \emph{¿La interfaz web cumple estándares de rendimiento,
accesibilidad, mejores prácticas y SEO?}

\textbf{Métrica:} categorías de Lighthouse \cite{lighthouse} (móvil, throttling Slow 4G)
sobre el login del frontend Angular.

\textbf{Datos crudos:} \texttt{docs/mediciones/lighthouse/lhci-20260730-2115.json} y
\texttt{...2117.json} (commit \texttt{1a07dc7}), más 6 corridas CI contra la URL
pública: escritorio d1--d3 (\texttt{lh-desktop-{1,2,3}.json}) y móvil m1--m3
(\texttt{lh-mobile-{1,2,3}.json}).

\textbf{Resultado:} Promedio de 6 corridas CI contra la URL pública: escritorio
Performance 95 / Accessibility 91 / Best Practices 92 / SEO 90; móvil Performance
77 / Accessibility 91 / Best Practices 92 / SEO 90. Los cuatro umbrales del
perfil de escritorio (80/90/90/90, \texttt{lighthouserc.js}) se cumplen. En móvil,
Performance (77) queda por debajo del umbral de 80, lo cual es esperable dado el
overhead del SPA en dispositivos móviles.

\textbf{Interpretación:} el SPA cumple los estándares en escritorio; en móvil,
Performance (77) es el punto más cercano al umbral y se documenta como mejora
menor (metadatos, optimización de bundle).

\section{Síntesis}
\label{sec:sintesis}

Table~\ref{tab:sintesis} resume el cumplimiento por bloque.

\begin{table}[H]
\centering
\caption{Results synthesis by evaluation block.}
\label{tab:sintesis}
\begin{tabular}{@{}llll@{}}
\toprule
Bloque & Indicador & Valor & Cumple \\
\midrule
C.1 Rendimiento & p95 $<$ 200\,ms (local, sin throttling) & 173.13\,ms & Sí \\
C.1 Rendimiento & p95 $<$ 200\,ms (Render, caché caliente, mismo usuario) & 504.77\,ms & No \\
C.1 Rendimiento & error rate $<$ 1\,\% & 0\,\% & Sí \\
C.2 Seguridad & SQL dinámico / A01--A09 & 0 violaciones & Sí \\
C.3 Usabilidad & SUS $\ge$ 70 & 63.0 & Marginal \\
C.4 Cobertura & líneas $\ge$ 70\,\% & 96.19\,\% & Sí \\
C.5 Calidad web & Perf/Acc/BP/SEO $\ge$ 80/90/90/90 & 95/91/92/90 & Sí \\
\bottomrule
\end{tabular}
\end{table}